\uuid{1V7Z}
\titre{Changements de variables}
\auteur{hanani ; mzari}
\organisation{exo7}
\licence{CC BY-NC-SA 4.0}
\datecreate{2020-01-20}
\niveau{L1}
\module{Analyse}
\chapitre{Calcul d'intégrales}
\sousChapitre{Autre}
\difficulte{2}
\theme{Calculs d'intégrales, changements de variables}
\exo7qcmid{lille-2.integrales.013}

\contenu{
\begin{question}
\qclassification{127.99}
\qtags{motcle=changements de variables}
Parmi les affirmations suivantes, cocher celles qui sont vraies :
\begin{answers}
\good{Le changement de variable $t=\sin x$ donne $\displaystyle \int _{\pi/6}^{\pi/4}\frac{\mathrm{d}\, x}{\sin x\tan x}=\int _{1/2}^{1/\sqrt{2}}\frac{\mathrm{d}\, t}{t^2}$.}
\bad{$\displaystyle \int _{\pi/6}^{\pi/4}\frac{\mathrm{d}\, x}{\sin x\tan x}=\frac{1}{\sqrt{2}}-\frac{1}{2}$.}
\bad{Le changement de variable $t=\cos x$ donne $\displaystyle \int _{0}^{\pi/3}\sin x\mathrm{e}^{\cos x}\mathrm{d}\, x=\int _1^{1/2}\mathrm{e}^t\mathrm{d}\, t$.}
\good{$\displaystyle \int _{0}^{\pi/3}\sin x\mathrm{e}^{\cos x}\mathrm{d}\, x=\mathrm{e}-\sqrt{\mathrm{e}}$.}
\end{answers}
\vskip3mm
\begin{explanations}
Avec $t=\sin x\Rightarrow \mathrm{d}t=\cos x\, \mathrm{d}x$, on a : $t(\pi/6)=1/2$, $t(\pi/4)=1/\sqrt{2}$ et
$$\displaystyle \int _{\pi/6}^{\pi/4}\frac{\mathrm{d}\, x}{\sin x\tan x}=\int _{\pi/6}^{\pi/4}\frac{\cos x\mathrm{d}\, x}{\sin ^2x}=\int _{1/2}^{1/\sqrt{2}}\frac{\mathrm{d}\, t}{t^2}=\left[-\frac{1}{t}\right]_{1/2}^{1/\sqrt{2}}=2-\sqrt{2}.$$
Avec $t=\cos x\Rightarrow \mathrm{d}t=-\sin x\, \mathrm{d}x$, on a : $t(0)=1$, $t(\pi/3)=1/2$ et
$$\displaystyle \int _{0}^{\pi/3}\sin x\mathrm{e}^{\cos x}\mathrm{d}\, x=\int _{1/2}^1\mathrm{e}^t\mathrm{d}\, t=\mathrm{e}-\sqrt{\mathrm{e}}.$$
\end{explanations}
\end{question}
}
